% Drawing 77/05-2 -- folder 77, batch 05, archive page 83 (PDF page 84).
% Transcription: dans la marge gauche, deux droites horizontales marquées $\Delta=0$ coupées par une courbe marquée $2=0$.
% Kind: two lines Delta = 0 cut by a curve 2 = 0. Method 2 (schematic redraw from the facsimile). Status: ok.
% Notes: Two nearly horizontal lines cut by a curve that crosses them; small cross marks on the curve; 2 = 0 at the top of the curve (and again at its foot). The Delta = 0 labels are at the right ends of the lines, just outside the crop: placed at the edge.
% Coordinates: 1 unit = 1 mm = 1/100 of the crop width (bbox in page fractions: [0.07, 0.28, 0.25, 0.48]).
% Redrawn by Claude (claude-opus-5-5), issue #73, 2026-11-07; not checked by a human; not compiled (no TeX in the session).
% Rights: reproduces part of the work -- behind RENDER_TRANSCRIPTS, never published by itself.
\documentclass[tikz,border=3pt]{standalone}
\usepackage{amsmath,amssymb}
\usetikzlibrary{patterns,arrows.meta}
\tikzset{pencil/.style={black!45}, struck/.style={black!35}, mark/.style={circle,draw,inner sep=1.3pt}, pt/.style={circle,fill,inner sep=0.9pt}, opt/.style={circle,draw,fill=white,inner sep=1.1pt}}
\begin{document}
\begin{tikzpicture}[x=1mm,y=1mm,line cap=round,line join=round,font=\scriptsize]
\draw plot[smooth, tension=0.55] coordinates {(40,-52) (70,-46) (100,-48)};
\draw plot[smooth, tension=0.55] coordinates {(38,-98) (70,-92) (100,-95)};
\draw plot[smooth, tension=0.55] coordinates {(82,-30) (72,-60) (72,-80) (78,-100) (82,-115) (72,-155)};
\draw (65,-58) -- (70,-72);
\draw (80,-108) -- (84,-122);
\node[inner sep=1pt] at (76,-20) {\tiny $2=0$};
\node[inner sep=1pt] at (92,-43) {\tiny $\Delta=0$};
\node[inner sep=1pt] at (92,-89) {\tiny $\Delta=0$};
\end{tikzpicture}
\end{document}
